% -----------------------------------------------------------------------------------
\section{How non-computable an optimal agent must be}\label{sec:degrees}

Theorem~\ref{thm:existence} and Theorems~\ref{thm:fixedcopy}--\ref{thm:exactsel} leave a precise
question: what oracle suffices, and what oracle is necessary, to compute an optimal decision
function? The natural guess is the halting problem, since admissibility of a deterministic
environment is a halting condition. The guess is wrong. Optimisation does not require recognising
admissibility, and the exact strength it requires is that of choosing, among
finitely many candidates, one that is never refuted. That strength is a classical degree class, the
PA degrees, and it rises to the halting problem and beyond only when the theory is also required to
say something definite off the admissible class or to break ties in a prescribed way.

Write $\le_T$ for Turing reducibility, $\mathbf 0'$ for the degree of the halting problem and
$\mathbf 0''$ for that of its relativisation. An oracle $X$ \emph{computes} an optimal agent of a
given totality if some $X$-computable agent of that totality is optimal.

\subsection{Diagonally non-computable functions}

\begin{definition}[Binary DNC and PA degree]\label{def:dnc}
A total function $g:\mathbb N\to\{0,1\}$ is \emph{binary diagonally non-computable} if
$g(e)\neq\varphi_e(e)$ whenever $\varphi_e(e)$ halts. A degree is a \emph{PA degree} if it computes
such a function.
\end{definition}

No binary DNC function is computable, since an index $e$ of it would give $g(e)\neq g(e)$. The name
comes from the equivalent condition that the degree computes a complete consistent extension of
Peano arithmetic. Equivalently, a degree is PA iff it computes a set separating $K_0$ from $K_1$, since
$1-s$ is binary DNC when $s$ separates them, and conversely. These equivalences go back to
\textcite{scott1962}, Solovay \autocite[Thm.~V.5.36]{odifreddi1989}, and
\textcite{jockuschsoare1972trans}; \textcite[§3.1]{hirschfeldt2014} gives a
modern account. The one consequence needed here is elementary.

\begin{lemma}[Finite selection]\label{lem:finitesel}
Let a finite nonempty set $B$ of candidates be given, together with a computably enumerable process
excluding candidates, uniformly in the data, such that at least one candidate is never excluded. Any
binary DNC function selects, with finitely many queries and uniformly in the data, a candidate that is
never excluded.
\end{lemma}

\begin{proof}
If $|B|>1$, split $B$ into nonempty parts $B_0,B_1$, and compute an index $k$ of the program that
ignores its input, runs the exclusion process, and outputs $i$ as soon as every member of $B_i$ has
been excluded. At most one part is ever wholly excluded, because some candidate survives. Continue
with $B_{g(k)}$: if it were wholly excluded, $\varphi_k(k)=g(k)$, contrary to the definition of $g$. So
$B_{g(k)}$ contains a survivor, and after fewer than $|B|$ steps a singleton survivor remains.
\end{proof}

The procedure never waits for the exclusion process; it only asks $g$ about programs that would wait.
Selecting a survivor is therefore weaker than recognising one; in the language of
\textcite{brattkaetal2012} it is finite closed choice.

\subsection{The degree theorem}

\begin{theorem}[Optimality has PA strength]\label{thm:pa}
For an oracle $X$:
\begin{enumerate}[label=(\roman*)]
  \item $X$ computes a globally total optimal agent on $\mathcal P_{\mathbb Q}$ iff $X$ has PA degree;
  \item $X$ computes a globally total optimal agent on $\mathcal P_{\mathrm{as}}$ iff $X$ has PA degree;
  \item $X$ computes a promise-total optimal agent on $\mathcal P_{\mathrm{as}}$ iff $X$ has PA degree;
\end{enumerate}
while a promise-total optimal agent on $\mathcal P_{\mathbb Q}$ is computable
(Theorem~\ref{thm:promiseopt}).
\end{theorem}

\begin{proof}
\emph{Upper bound for (i).} On an owned query of a description $E$, take the policies of $E$ as
candidates, and let the exclusion process simulate $E^\pi$ for every $\pi$; if all simulations halt,
it excludes exactly the non-maximising policies, and otherwise it excludes nothing. Some candidate is
never excluded on every description, and on admissible $E$ the survivors are exactly the maximisers.
Lemma~\ref{lem:finitesel} selects one with an $X$-computable binary DNC function; every owned query of
$E$ reconstructs the same selection, and the agent answers with its component. It halts on every
input, since the lemma makes finitely many queries and never waits.

\emph{Upper bound for (ii) and (iii).} For a policy $\pi$ and a time $s$, let $H_{\pi,s}$ be the
probability that $E^\pi$ has halted by step $s$ and $L_{\pi,s}$ the expected reward collected on those
executions. Both are rationals computable from the finite execution tree, and with
$B_{\pi,s}=L_{\pi,s}+1-H_{\pi,s}$ the intervals $[L_{\pi,s},B_{\pi,s}]$ are nested, because newly
halted mass contributes a reward between $0$ and that mass. On admissible $E$,
\begin{equation}\label{eq:intervals}
  V_E(\pi)\in[L_{\pi,s},B_{\pi,s}],\qquad B_{\pi,s}-L_{\pi,s}=1-H_{\pi,s}\to0 .
\end{equation}
Exclude $\pi$ as soon as $L_{\rho,s}>B_{\pi,s}$ for some $\rho$ and $s$. On admissible environments
this excludes exactly the strictly suboptimal policies. On every description some policy survives:
let $\ell_\pi=\lim_sL_{\pi,s}$, the reward collected on terminating executions; then
$L_{\pi,s}\le\ell_\pi\le B_{\pi,s}$, so a maximiser $\pi^\circ$ of $\ell$ cannot be excluded, since
$\ell_\rho\ge L_{\rho,s}>B_{\pi^\circ,s}\ge\ell_{\pi^\circ}$ would follow. The quantity $\ell$ serves only
this argument and changes no admissibility condition. Lemma~\ref{lem:finitesel} now gives a globally
total agent, optimal on $\mathcal P_{\mathrm{as}}$, and \emph{a fortiori} a promise-total one.

\emph{Lower bound for (i).} Let $D_e$ run $\varphi_e(e)$ and, if it returns a bit $b$, pay
$\one[a=b]$; if it returns anything else, pay $0$ for both actions. $D_e$ is admissible iff
$\varphi_e(e)$ halts. For a globally total optimal $T$, $s(e)=T(d(D_e))$ is total, even where $D_e$ is
inadmissible, and equals $b$ whenever $e\in K_b$. So $s$ separates $K_0$ from $K_1$ and $1-s$ is binary
DNC.

\emph{Lower bound for (ii) and (iii).} Every environment $E_e$ of Theorem~\ref{thm:exactsel} lies in
$\mathcal P_{\mathrm{as}}$, and an optimal agent, promise-total or global, satisfies $T(d(E_e))=b$ for
$e\in K_b$. Since all $E_e$ are admissible, $s(e)=T(d(E_e))$ is total, and it separates $K_0$ from
$K_1$.
\end{proof}

Three consequences of the theorem bear on the thesis. First, the deterministic lower bound uses no
self-reference and no copy of any agent: $D_e$ is generated independently of the theory evaluated.
It contains Theorem~\ref{thm:fixedcopy} as the special case $X=\varnothing$, and shows that the
obstruction to a globally total optimiser is not a predictor but the task of completing the promise
optimiser of Theorem~\ref{thm:promiseopt} on descriptions that might be admissible. Second, the
optimiser does not decide admissibility, which for $\mathcal P_{\mathrm{as}}$ is of degree
$\mathbf 0''$ (Proposition~\ref{prop:admissibility}); it only selects a survivor. Third, on
$\mathcal P_{\mathrm{as}}$ the strength is required even of agents that halt only on the promise,
which is why Theorem~\ref{thm:exactsel} is a negative result about exact selection and not about
totality.

\begin{corollary}[The spectrum of optimal agents]\label{cor:spectrum}
On $\mathcal P_{\mathbb Q}$ and on $\mathcal P_{\mathrm{as}}$:
\begin{enumerate}[label=(\roman*)]
  \item every PA degree is the exact Turing degree of some optimal decision function;
  \item some optimal decision function $T$ satisfies $\mathbf 0<_T T<_T\mathbf 0'$;
  \item for every non-computable set $Y$ there is an optimal decision function that does not compute
        $Y$, so no single non-computable set is computed by every optimal decision function.
\end{enumerate}
\end{corollary}

\begin{proof}
(i) Given $X$ of PA degree, take an $X$-computable optimiser from Theorem~\ref{thm:pa} and modify it on
a computable family of distinct binary environments $F_n$ in which both actions tie, answering
$X(n)$ on $d(F_n)$. Ownership confines each modification to $F_n$, so the result is still optimal and
$X$-computable, and it computes $X$. (ii) By the low basis theorem \autocite[Thm.~2.1]{jockuschsoare1972trans}, which applies
to the $\Pi^0_1$ class of separating sets of $K_0$ and $K_1$,
there is a PA degree $X$ with $X'\equiv_T\mathbf 0'$; it does not compute $\mathbf 0'$, since otherwise
$\mathbf 0''\le_T X'\equiv_T\mathbf 0'$. Apply (i). (iii) By the cone-avoidance basis theorem for
$\Pi^0_1$ classes \autocite[Thm.~2.5]{jockuschsoare1972trans}, applied to the class of separating sets of $K_0$
and $K_1$, some PA degree does not compute $Y$. Apply (i).
\end{proof}

So \enquote{optimal agents need the halting problem} is false, and there is no least degree of
optimality: the PA degrees have no least element.

\subsection{Admissibility and prescribed answers}

\begin{proposition}[Complexity of admissibility]\label{prop:admissibility}
Membership in $\mathcal P_{\mathbb Q}$ is $\Sigma^0_1$-complete, and membership in
$\mathcal P_{\mathrm{as}}$ is $\Pi^0_2$-complete.
\end{proposition}

\begin{proof}
Since there are finitely many policies, $E\in\mathcal P_{\mathbb Q}$ iff for some $s$ every $E^\pi$
halts legally within $s$ steps, a decidable condition on $s$; an environment that ignores the action,
runs a given computation and then pays $0$ is admissible exactly when the computation halts. For the
stochastic class, $E\in\mathcal P_{\mathrm{as}}$ iff $\forall k\,\exists s\,\forall\pi\in\Pol_E\colon
H_{\pi,s}>1-2^{-k}$, and the last quantifier is finite. For hardness, reduce
$\mathrm{TOT}=\{e:\varphi_e\text{ is total}\}$: draw a geometric $N$ and run $\varphi_e(N)$ without a
bound. Every $n$ has positive probability, so the environment terminates almost surely exactly when
$e\in\mathrm{TOT}$.
\end{proof}

The stochastic half is the $\Pi^0_2$-completeness of almost-sure termination proved by
\textcite{kaminskikatoen2015}; see also \textcite[§4]{kaminskietal2019}. A theory that must answer a prescribed default on every inadmissible
description, as the existence theorem permits, needs this much more.

\begin{proposition}[Prescribed answers]\label{prop:defaults}
Let the agent be required to return a given computable default on every description that is not an
owned query of an admissible environment.
\begin{enumerate}[label=(\roman*)]
  \item On $\mathcal P_{\mathbb Q}$ every such optimal agent computes $\mathbf 0'$, and one is computable
        from $\mathbf 0'$.
  \item On $\mathcal P_{\mathrm{as}}$ every such optimal agent computes $\mathbf 0''$, and one is
        computable from $\mathbf 0''$.
  \item On $\mathcal P_{\mathrm{as}}$, if ties must be broken in a fixed order, every promise-total
        optimal agent computes $\mathbf 0'$, and one is computable from $\mathbf 0'$.
\end{enumerate}
\end{proposition}

\begin{proof}
Take the default to be $0$; for another computable default, the environments below compute it on their
own description and reward its opposite. (i) Let $H_e$ wait for $\varphi_e(e)$ and, if it halts, pay
the action $a\in\{0,1\}$. If $\varphi_e(e)$ halts, $1$ is the unique optimum; otherwise $H_e$ is
inadmissible and the default applies. So $T(d(H_e))=1$ iff $\varphi_e(e)$ halts. Conversely,
$\mathbf 0'$ decides admissibility, after which the simulation of Theorem~\ref{thm:promiseopt}
applies. (ii) Let the environment draw a geometric $N$, run $\varphi_e(N)$, and if it halts pay $a$. It
is admissible, with unique optimum $1$, iff $e\in\mathrm{TOT}$, so $T$ decides $\mathrm{TOT}$.
Conversely, $\mathbf 0''$ decides admissibility, and on admissible environments $\mathbf 0'$ decides
the $\Sigma^0_1$ exclusion predicate of the proof of Theorem~\ref{thm:pa}, whose survivors are the
maximisers. (iii) Let action $0$ pay $0$, and let action $1$ draw a geometric $N$, simulate a given
computation for $N$ steps and pay $1$ if it has halted. Action $1$ has value $2^{-t}$ if the
computation halts at step $t$ and value $0$ otherwise, so with ties broken towards $0$ the agent
chooses $1$ iff the computation halts. Conversely, $\mathbf 0'$ decides exclusion, and the first
survivor in the fixed order is the prescribed maximiser.
\end{proof}

The PA classification therefore concerns optimality as defined, which permits any maximiser and any
answer off the admissible class. Every additional specification of what the theory must say, such as
a default outside the class or an order among tied alternatives, specifies something other than the
choice among alternatives in admissible problems, and it is what raises the degree.

\subsection{Relativisation}

\begin{theorem}[Relativised diagonal]\label{thm:reldiag}
Fix an oracle $X$, and let environments and their embedded programs be $X$-computable, with the same
ownership and termination conditions. Then no globally total $X$-computable agent is optimal, an
optimal decision function still exists, and Theorem~\ref{thm:pa} and Proposition~\ref{prop:defaults}
hold with binary $\mathrm{DNC}^X$ functions in place of binary DNC functions and with $X'$, $X''$ in
place of $\mathbf 0'$, $\mathbf 0''$.
\end{theorem}

\begin{proof}
The relativised recursion theorem \autocite[Thm.~II.2.13 and p.~177]{odifreddi1989} constructs, for each globally total
$X$-computable $T$, an $X$-computable environment with $U(E^X_T,A)=\one[A(d)\ne T(d)]$; its fixed
computation of $T(d)$ halts, so it is admissible and defeats $T$. Existence uses only finiteness and
ownership. The remaining proofs relativise line by line.
\end{proof}

The oracle of the environments is held fixed in the theorem. If environments may consult the
optimiser's stronger oracle, the class is a new one and the diagonal applies again: for every class of
environments closed under an oracle, an optimal agent exists, and every optimal agent that halts on
every input lies strictly above that oracle in the sense of computing a $\mathrm{DNC}^X$ function. This
makes precise the claim that an optimal theory must lie above every environment that can run it.

\subsection{Sharpness of the hypotheses}

Table~\ref{tab:sharpness} records what fails when each hypothesis is dropped. The constructions are
short, and the non-obvious ones are given after the table.

\begin{table}[htbp]
  \centering\small
  \renewcommand{\arraystretch}{1.2}
  \captionabove{Hypotheses of Theorems~\ref{thm:existence} and~\ref{thm:pa}, and what fails without
                them.}
  \label{tab:sharpness}
  \begin{tabularx}{\textwidth}{@{}>{\raggedright\arraybackslash}p{9.6em}
                                    >{\raggedright\arraybackslash}X@{}}
    \toprule
    hypothesis dropped & consequence \\
    \midrule
    owned queries & no optimal function of any degree (Prop.~\ref{prop:foreign}) \\
    embedded copies fixed & the diagonal of Thm.~\ref{thm:fixedcopy} disappears; the lower bounds of
                            Thm.~\ref{thm:pa} remain, since they use no copy \\
    termination under every table & with divergence paying a fixed value, exact optimisation decides
                                    halting even on deterministic environments \\
    explicit rational values & with values given as computable reals, deterministic exact
                               optimisation needs PA degree even on the promise \\
    rewards bounded in $[0,1]$ & an almost surely terminating family with finite values forces
                                 $\mathbf 0'$, even for $\varepsilon$-optimality with $\varepsilon<\tfrac12$ \\
    exact optimality & $\varepsilon$-optimal and unique-maximiser agents are computable on the
                       promise (Prop.~\ref{prop:approxopt}) \\
    finite policy spaces & a maximum need not exist: payoffs $1-1/(n+1)$ on actions $n\in\mathbb N$ \\
    \bottomrule
  \end{tabularx}
\end{table}

For computable reals, let $v_e(b)=2^{-t}$ and $v_e(1-b)=0$ if $\varphi_e(e)$ first halts with bit $b$
at step $t$, and $v_e\equiv0$ otherwise. These values have uniformly computable fast Cauchy names, and
every one of them is in fact rational; what is missing is its explicit exact representation. An exact
optimiser separates $K_0$ from $K_1$ as in Theorem~\ref{thm:pa}. For unbounded rewards, draw $N\ge1$
with probability $2^{-N}$; let action $0$ pay $\tfrac12$ and action $1$ pay $2^N$ exactly when a given
computation first halts at step $N$. Then $V(1)=1$ if the computation halts and $V(1)=0$ otherwise, so
any agent within $\tfrac12$ of optimal decides halting. The values are finite but not uniformly
computable, which is why Proposition~\ref{prop:approxopt}(i) does not apply.

\subsection{Unfair problems}\label{sec:unfair}

The remaining question is the one the thesis began with: which theories face problems built against
them? \enquote{Built against} names a provenance unless a mathematical witness is specified. The
anti-prediction environments $E_p$ of Proposition~\ref{prop:antipred} supply one without reference to
optimality: $E_p$ is \emph{accurate for} a decision function $T$ if $p(d(E_p))=T(d(E_p))$. For
admissible $E_p$, accuracy is equivalent to defeat, since $T$ then receives $0$ while the optimum is
$1$.

\begin{proposition}[Unfairness and computability]\label{prop:unfair}
\leavevmode
\begin{enumerate}[label=(\roman*)]
  \item For every oracle $X$, every globally total $X$-computable decision function has an accurate,
        admissible, $X$-computable anti-prediction environment, constructed uniformly from its index.
  \item No optimal decision function has an accurate admissible anti-prediction environment.
  \item Some non-computable decision function has an accurate admissible anti-prediction environment.
\end{enumerate}
\end{proposition}

\begin{proof}
(i) is Theorem~\ref{thm:reldiag} and (ii) is Proposition~\ref{prop:antipred}. For (iii), let $T^*$ be
optimal and non-computable, by Theorem~\ref{thm:pa} and Corollary~\ref{cor:spectrum}, and let $E_0$ be
the anti-prediction environment of the constant program $0$, so that $T^*(d(E_0))=1$. Changing this
single answer to $0$ gives a non-computable function for which $E_0$ is accurate.
\end{proof}

The proposition refutes the claim that unfair problems exist exactly for computable theories. The
defensible direction is
\[
  \text{global simulability by the environments}\ \Longrightarrow\ \text{uniformly constructible
  defeat},
\]
and neither defeat nor local accuracy implies computability: non-computable theories can make
predictable mistakes. Nor can computability be read off the restriction to owned queries, since every
such restriction is a finite table and hence computable. Within the anti-prediction family the
equivalence of accuracy and defeat is exact, but it characterises the family, not the computable
theories. \textcite[Thm.~18]{leikehutter2015} make a related observation in a different setting:
finite hard-coded behaviour defines a computable environment even when the policy it is aimed at is
not computable.

\subsection{Resource bounds}

Finite clocks or circuit descriptions can support globally total optimisers, because every required
evaluation then has a finite bound: on a decidable clocked language with a supplied execution bound,
deterministic values are evaluated exhaustively and stochastic ones by summing a finite rational
execution tree. They do not make optimisation as cheap as one execution.

\begin{proposition}\label{prop:sat}
If an agent running in time polynomial in the length of the description is exactly optimal on all
environments with $n$ binary information sets and a Boolean circuit as payoff, then
$\mathrm{P}=\mathrm{NP}$. If it is exactly optimal on all single-choice environments that run in
polynomial time with fair coins, then $\mathrm{P}=\mathrm{PP}$.
\end{proposition}

\begin{proof}
Given a Boolean formula with $n$ variables, form the environment whose information sets are the
variables and whose payoff is $1$ exactly on satisfying assignments. Query the agent on its $n$ owned
queries and evaluate the formula on the resulting policy. By optimality the policy is satisfying
whenever a satisfying assignment exists, so this decides satisfiability in polynomial time
\autocite{cook1971}. For the second claim, let action $1$ evaluate the formula on a uniformly random
assignment and pay $1$ if it is satisfied, and let action $0$ pay $1$ with probability
$\tfrac12+2^{-(n+1)}$. The value of action $1$ is a multiple of $2^{-n}$ and that of action $0$ is
not, so there are no ties, and action $1$ is optimal iff more than half of the assignments satisfy the
formula. This decides the $\mathrm{PP}$-complete problem $\mathrm{MAJSAT}$ \autocite{gill1977}.
\end{proof}

Each play of these environments is short; the difficulty lies in comparing alternatives, all of
whose executions must be accounted for. A resource-bounded form of Theorem~\ref{thm:fixedcopy} holds
against an agent $T$ whenever the environments' budget accommodates self-description, a simulation of
$T$ on that description, and the final comparison. If the budget cannot accommodate the simulation,
the diagonal gives no lower bound, so the useful resource condition is that the environments can run
the agent, not a ranking of agents and environments by a single bound.

Per-input runtime bounds do not decide program equivalence either; the construction of
Theorem~\ref{thm:haltred} uses components bounded by $|x|$ steps on input $x$. The bounded variants
of Löb's theorem proved by \textcite{critch2019} yield positive cooperation results for specified
proof-responsive programs under proof-length assumptions; they do not establish optimality on a
class of environments of the present kind.

\subsection{What a fixed semantics settles}\label{sec:fixedsemantics}

Theorem~\ref{thm:promiseopt} extends at once to any environment class that supplies a
\emph{reference-independent} effective evaluator: a computable $V(E,\pi)$ on a finite policy space,
with exactly comparable values, whether the evaluator works by source analysis, by simulation, or by
proof. Finite circuits, explicitly clocked programs with stipulated time-out behaviour, and
descriptions carrying checkable certificates of the required evaluations all have such evaluators,
and admit optimisers that are globally total on the certified language, returning a default on
invalid strings. One might conclude that once a rule for propagating the dependence of the world on
the agent's output is fixed, optimal choice is a finite computation, and the only remaining question
is which rule.

The conclusion is false for anchored kernels. A fixed, computable kernel makes each comparison
computable, but it need not admit an agent whose actual output maximises it. In
Proposition~\ref{prop:doublecall} under $O_{\mathrm{const}}$ every baseline yields a finite,
unambiguous ranking, and every baseline is ranked below its alternative. Maximising a supplied table
$V_E$ is a finite computation; being a fixed point of a baseline-dependent kernel is a separate
problem, the ratifiability problem of Part~\ref{part:single} in program form.

Two further remarks fix the scope of \enquote{largest solvable class}. For reference-independent
problems with disjoint owned queries, every decidable class on which some globally total agent is
optimal can be enlarged by a single further problem, by altering the agent on that problem's owned
queries only; under $\cls_{\mathrm{self}}$ this changes nothing on the original class, so there is
no largest such class. Conversely, the class of all terminating, rational-valued, finite-policy
templates is a maximal promise domain for that representation, so an unqualified claim that there is
no largest solvable class would also be false.

Sections~\ref{sec:computation} and~\ref{sec:degrees} thus answer the question for a designated
semantics completely. Optimal choice exists problem by problem; an optimal function is
non-computable exactly when it must halt everywhere or compare computable reals exactly, and then its
degree is PA; every further demand on it raises the degree without touching the choice among
alternatives. Everything in this answer presupposed that the substitution interface was given. The
remaining sections remove it.
